\section{Evaluation}
\subsection{Exploratory SAFE baseline rerun}
We reran SAFE on the bundled Twitter and Wiki splits using one seed
and TensorFlow 2.12 in compatibility mode. The cumulative-hazard loss
was evaluated with an algebraically equivalent stable expression
after the original expression overflowed on Twitter. Model updates
used training data. Checkpoint and threshold selection maximized mean
validation accuracy over prefixes 1--5. Validation and test inference
received only the first five feature vectors per account. All test
records were scored, including remainder batches.

\paragraph{Split audit.}
We preserved the test splits and removed records from earlier splits
when their full feature sequences matched another split, giving test
priority over validation and validation priority over training.
On Twitter, this removed 64 training and 15 validation records,
leaving 3,806 training, 541 validation, and 1,114 test records.
Wiki had no exact full-sequence matches across splits and retained
1,231 training, 175 validation, and 353 test records. After filtering,
neither dataset had exact full-sequence matches across splits.
Exact matches are a feature-level check. Sparse sequences may match
across different accounts, and the bundled arrays lack account IDs and
source timestamps for verifying account separation and feature availability.

\begin{table}[t]
\centering
\small
\setlength{\tabcolsep}{4pt}
\begin{tabular}{llrrrr}
\hline
Dataset & Run & Prec. & Recall & F1 & Acc. \\
\hline
Twitter & Published & 82.0 & 55.7 & 65.4 & 71.8 \\
        & Audited rerun & 71.8 & 60.0 & 64.8 & 67.9 \\
Wiki    & Published & 71.1 & 88.0 & 78.7 & 76.4 \\
        & Audited rerun & 69.7 & 89.0 & 78.2 & 75.5 \\
\hline
\end{tabular}
\caption{SAFE performance (\%) averaged over prefixes 1--5.
The published study reports ten runs. The audited rerun uses one
seed per dataset.}
\label{tab:safe-audited-rerun}
\end{table}

\paragraph{Results and limitations.}
The published values are from Table~1 of \cite{zheng2019safe}.
In Table~\ref{tab:safe-audited-rerun}, local F1 was within 0.6
percentage points of the published score on both datasets. Twitter
precision remained 10.2 points below the published value.
The rerun changed overlap handling, batching, and the validation
window together. The legacy test results had been inspected before
these changes, so the rerun is exploratory. A confirmatory evaluation
requires a fresh account-disjoint holdout, verifiable feature
timestamps, and a protocol fixed before test evaluation.

\subsection{Planned SAFEPlus evaluation}
\paragraph{Stage 1: likelihood verification.}
Enumerate short sequences to compare Equations~\ref{eq:discrete-detected}
and~\ref{eq:discrete-censored} against the linear-time recurrence. Check
normalization, posterior sums, gradient calculations, variable sequence
lengths, and censoring boundaries. Generate an explicit detection observation
indicator separately from latent fraud status.

\paragraph{Stage 2: synthetic recovery.}
Generate known onset and delay distributions, then hide onset during fitting.
Vary delay length, censoring, sample size, and the fraction of audited onset
labels. Include the rate-swapping example as a deliberate failure case.
Compare no delay information, a known delay process, and partial onset
supervision across multiple initializations. Report held-out detection
log-likelihood separately from onset error, posterior interval coverage,
undetected-fraud calibration, and parameter sensitivity. Strong detection fit
with poor onset recovery is an important diagnostic outcome.

\paragraph{Stage 3: controlled comparison with SAFE.}
Establish the original data and loss conventions before porting the legacy
implementation to a current framework. Compare SAFE, a standard detection-time
survival model, and the proposed joint model under matched causal encoders,
splits, and tuning budgets. Add a supervised-onset oracle only in synthetic
experiments as a reference. Introduce an onset-dependent delay model only
after the restricted model is validated.

\paragraph{Real-data audit and metrics.}
Determine whether timestamps denote observations, edits, suspensions, or
intervals between status checks. Audit label semantics, censoring, padding,
and train/validation/test separation. Select alert thresholds on validation
data only. Report prefix-based precision--recall performance, false alerts,
and lead time relative to observed suspension, with repeated-run uncertainty.
Where chronological splits are possible, prevent future account or platform
information from entering training features. Without independently verified
onset labels, earlier alerts than suspension do not establish accurate fraud
onset recovery or anticipation before fraud occurs.
